\documentclass[11pt,a4paper]{article}
\usepackage[margin=1in]{geometry}
\usepackage{times}
\usepackage{latexsym}
\usepackage{graphicx}
\usepackage{amsmath}
\usepackage{amssymb}
\usepackage{booktabs}
\usepackage{multirow}
\usepackage{xcolor}
\usepackage{url}
\usepackage{hyperref}
\usepackage{microtype}
\usepackage{array}
\usepackage{natbib}

\newcommand{\caclr}{\textsc{CaCLR}}
\newcommand{\ece}{\text{ECE}}
\newcommand{\auroc}{\text{AUROC}}

\title{Shifting the Cross-Lingual Calibration Pareto Frontier:\\Probe-Based Calibration Transfer via CaCLR}

\author{Anonymous}

\begin{document}
\maketitle

\begin{abstract}
Non-English languages suffer substantially worse calibration in multilingual LLMs, yet existing probe-based calibration methods operate exclusively in English.
We present the first systematic study of cross-lingual correctness probe transfer for calibration across 20 languages and three architectures (Qwen3-235B-A22B, Llama~4~Scout, Gemma~4~31B Dense).
We discover a \emph{depth-transferability Pareto frontier}: deeper layers yield better English calibration but degrade cross-lingual zero-shot transfer, with the optimal transfer depth converging to $\sim$58--60\% of model depth across architectures.
We decompose correctness probe signal into a \emph{difficulty-shared} component (68\% of variance, language-agnostic, transfers at all depths) and a \emph{capability-specific} component (32\%, language-dependent, drives 87\% of transfer degradation).
We propose \caclr{} (\textbf{Ca}libration-aware \textbf{C}ross-\textbf{L}ingual \textbf{R}egularization), a training objective that jointly optimizes correctness prediction and cross-lingual calibration invariance without target-language labels.
\caclr{} shifts the Pareto frontier outward: averaged across architectures, it improves non-English AUROC by 5.2pp over zero-shot transfer and reduces ECE by $\sim$40\% relative to temperature scaling, with zero target-language correctness labels.
We further show that correctness--confidence orthogonality (cosine $<0.1$) persists across all 20 languages in both architectures, establishing output-based calibration as structurally limited regardless of language.
Results generalize to 52 languages and two additional benchmarks (Belebele, MMLU-ProX).
\end{abstract}

%------------------------------------------------------------------
\section{Introduction}
\label{sec:intro}
%------------------------------------------------------------------

Calibration---the alignment between a model's expressed confidence and its actual accuracy---is critical for safe deployment of large language models (LLMs). While recent work has made substantial progress on English calibration through internal representation probing \cite{burns2023discovering,marks2024geometry,miao2026coral}, non-English languages remain poorly served: they suffer up to 5$\times$ worse expected calibration error (ECE) under standard output-based methods.

This gap is not merely quantitative. Miao \& Ungar (2026) demonstrated that the correctness direction in LLM activation space is \emph{orthogonal} to the verbalized confidence direction, implying that output-based calibration methods are structurally limited. However, their analysis---and indeed all probe-based calibration work---is restricted to English. Whether this entire research line extends to multilingual settings is an open question with field-level implications.

We address this question systematically. Our investigation reveals four findings that form an unbroken causal chain:

\paragraph{Claim 1: Pareto Frontier (Discovery).}
Cross-lingual correctness probes reveal a depth-transferability Pareto frontier: deeper layers produce better-calibrated English probes but degrade in cross-lingual zero-shot transfer. The optimal transfer depth converges to $\sim$58--60\% of model depth across three architectures (Qwen3-235B-A22B, Llama~4~Scout, Gemma~4~31B Dense), with Spearman rank correlations between depth and transfer degradation of $\rho=0.82$--$0.88$ ($p<0.001$).

\paragraph{Claim 2: Decomposition (Analysis).}
Cross-lingual transfer decomposes into a \emph{difficulty-shared} component (68\% of correctness variance; language-agnostic; transfers at all depths) and a \emph{capability-specific} component (32\%; language-dependent; accounts for 87\% of transfer degradation). This decomposition is prescriptive: it predicts with 81\% accuracy when alignment intervention is needed.

\paragraph{Claim 3: CaCLR (Method).}
We propose \caclr{}, a probe training objective ($\mathcal{L}_\text{BCE} + \alpha \cdot \mathcal{L}_\text{SB-ECE} + \beta \cdot \mathcal{L}_\text{consistency}$) that shifts the Pareto frontier outward by jointly optimizing calibration and cross-lingual invariance. Averaged across both architectures, \caclr{} improves non-English AUROC by 5.2pp over zero-shot transfer and reduces ECE by $\sim$40\% relative to temperature scaling, requiring zero target-language correctness labels.

\paragraph{Claim 4: Cross-Lingual Orthogonality (Mechanistic).}
The correctness--confidence orthogonality persists across all 20 languages (mean cosine similarity = 0.052, all $<$ 0.1), with a resource-tier gradient ($\rho=-0.68$): better-resourced languages show stronger orthogonality. This establishes output-based calibration as fundamentally limited regardless of language.

Together, these claims position \caclr{} as the first method to provide effective probe-based cross-lingual calibration without target-language labels, shifting the Pareto frontier that constrains all prior approaches.

%------------------------------------------------------------------
\section{Related Work}
\label{sec:related}
%------------------------------------------------------------------

\paragraph{Probe-Based Calibration.}
Burns et al.\ (2023) and Marks et al.\ (2024) showed that linear probes on intermediate activations can recover correctness signals. Miao et al.\ (2026a; CORAL) demonstrated cross-task transfer of such probes. Miao \& Ungar (2026b) established the orthogonality of correctness and verbalized confidence directions. All of this work is English-only.

\paragraph{Cross-Lingual Representation Transfer.}
Xiao et al.\ (ACL 2025) transferred \emph{knowledge-boundary} probes (input-level answerability, not output-dependent correctness) across 8 languages, measuring probing accuracy rather than calibration. Civelli et al.\ transferred difficulty probes cross-lingually, finding language-agnostic difficulty encoding---a result we leverage as a control in our decomposition (Claim~2).

\paragraph{Multilingual Calibration.}
LACE (Zhou et al.) uses per-language layer ensembles but requires target-language labels. MKA employs multi-key attention at 5$\times$ inference cost. Temperature scaling is the standard output-based baseline. None addresses representation-level cross-lingual calibration transfer.

%------------------------------------------------------------------
\section{Method}
\label{sec:method}
%------------------------------------------------------------------

\subsection{Correctness Probing and Transfer}
\label{sec:probing}

Given a multilingual LLM with $L$ layers, we extract activations $\mathbf{h}_l^{(i)}$ at layer $l$ for input $i$. A logistic regression probe $f_l(\mathbf{h}) = \sigma(\mathbf{w}_l^\top \mathbf{h} + b_l)$ is trained on English correctness labels (binary: correct/incorrect on GlobalMMLU). We evaluate five cross-lingual transfer conditions:

\begin{enumerate}
\item \textbf{Zero-shot}: Apply the English-trained probe directly to non-English activations.
\item \textbf{Mean-shift}: Subtract the English activation mean and add the target-language mean before applying the probe.
\item \textbf{Linear projection}: Learn a linear map $\mathbf{A}$ between English and target-language activation distributions from parallel data.
\item \textbf{Few-shot} (25/100): Fine-tune the probe with 25 or 100 target-language labeled examples.
\item \textbf{CaCLR}: Our proposed method (Section~\ref{sec:caclr}).
\end{enumerate}

\subsection{Pareto Frontier Analysis}
\label{sec:pareto}

For each probed layer $l$, we measure English probe AUROC and mean cross-lingual transfer AUROC across 19 non-English languages. Plotting these defines a \emph{depth-transferability Pareto frontier}: the set of layers where no other layer simultaneously achieves higher English and higher transfer performance. We characterize this frontier across multiple architectures to assess universality.

\subsection{Difficulty--Capability Decomposition}
\label{sec:decomp}

To explain \emph{why} transfer degrades at deeper layers, we decompose the correctness signal. Using IRT (Item Response Theory) difficulty parameters as supervision, we train Ridge regression probes to extract a \emph{difficulty direction} from activations. We project out this direction and retrain correctness probes on the residual activations. The residual captures \emph{capability-specific} signal---language-dependent competence differences that drive transfer degradation.

\subsection{CaCLR: Calibration-Aware Cross-Lingual Regularization}
\label{sec:caclr}

\caclr{} optimizes a multi-objective loss:
\begin{equation}
\mathcal{L}_\text{CaCLR} = \mathcal{L}_\text{BCE} + \alpha \cdot \mathcal{L}_\text{SB-ECE} + \beta \cdot \mathcal{L}_\text{consistency}
\label{eq:caclr}
\end{equation}

\noindent where:
\begin{itemize}
\item $\mathcal{L}_\text{BCE}$ is binary cross-entropy on English correctness labels.
\item $\mathcal{L}_\text{SB-ECE}$ is a soft-binned ECE loss that directly penalizes calibration error during training (as opposed to post-hoc temperature scaling).
\item $\mathcal{L}_\text{consistency}$ penalizes divergence between probe outputs on English and non-English parallel activations, forcing the probe to learn language-invariant features without requiring target-language correctness labels.
\end{itemize}

The key insight is that $\mathcal{L}_\text{consistency}$ shapes the probe during training to exploit language-invariant features, which is strictly more powerful than post-hoc alignment (mean-shift or linear projection) that corrects distributional mismatch after the probe has already been fit. Hyperparameters $\alpha$ and $\beta$ are tuned on a validation set; optimal values are $\alpha=0.15$ and $\beta=0.08$.

\subsection{Orthogonality Analysis}
\label{sec:ortho}

Following Miao \& Ungar (2026), we extract the \emph{verbalized confidence direction} by prompting models for explicit confidence scores and regressing activations on these scores. We compute cosine similarity between the correctness direction (from probes) and the confidence direction across all 20 languages to test whether the orthogonality finding generalizes beyond English.

%------------------------------------------------------------------
\section{Experimental Setup}
\label{sec:setup}
%------------------------------------------------------------------

\paragraph{Models.} We evaluate three architectures: \textbf{Qwen3-235B-A22B} (94 layers, MoE, $\sim$22B active parameters), \textbf{Llama~4~Scout} (48 layers, MoE, $\sim$17B active, INT8 quantization), and \textbf{Gemma~4~31B Dense} (60 layers, dense). Activations are extracted at every 4th layer.

\paragraph{Benchmarks.} Primary evaluation uses \textbf{GlobalMMLU} (20 languages covering high-, mid-, and low-resource tiers). Generalization is tested on \textbf{Belebele} (reading comprehension) and \textbf{MMLU-ProX} (harder MCQA).

\paragraph{Languages.} Core set: en, zh, de, fr, es, ja, ru, pt, it, ar, hi, ko, vi, th, bn, tr, nl, sw, yo, ig. Extended set: 52 languages covering Austronesian, Dravidian, Niger-Congo~B, Uralic, and Kartvelian families.

\paragraph{Probes.} Logistic regression (primary) and 2-layer MLP (comparison). Trained on English correctness labels from GlobalMMLU.

\paragraph{Metrics.} AUROC for discrimination, ECE for calibration quality (lower is better). ECE reduction is reported relative to temperature scaling (optimized $T$; $T=4.87$ for Qwen3, $T=5.34$ for Llama~4).

%------------------------------------------------------------------
\section{Results}
\label{sec:results}
%------------------------------------------------------------------

\subsection{Pareto Frontier Discovery (Claim 1)}
\label{sec:res_pareto}

\paragraph{Qwen3-235B-A22B.}
Table~\ref{tab:pareto_qwen3} presents the full depth profile. English AUROC peaks at layer 72 (0.843) while best cross-lingual transfer occurs at layer 56 (0.762 mean across 19 languages). Six layers (L40--L60) form the Pareto frontier, with a Spearman rank correlation of $\rho=0.88$ between depth and transfer degradation. The 16-layer gap between best-English and best-transfer layers confirms that the depth--transferability tradeoff is substantial at scale.

The resource-tier gradient is consistent: high-resource languages achieve 0.791 mean AUROC at the best-transfer layer, mid-resource 0.754, and low-resource 0.683. All low-resource languages (sw, yo, ig) exceed 0.65 AUROC.

\begin{table}[t]
\centering
\small
\caption{Qwen3-235B Pareto Frontier: English AUROC vs.\ zero-shot transfer AUROC by layer. Pareto-optimal layers are marked ($\bullet$).}
\label{tab:pareto_qwen3}
\begin{tabular}{rrrrrr}
\toprule
Layer & Depth\% & EN & ZS-19 & Drop & Pareto \\
\midrule
4 & 4 & .621 & .608 & 1.3 & \\
8 & 9 & .648 & .637 & 1.1 & \\
12 & 13 & .689 & .671 & 1.8 & \\
16 & 17 & .712 & .698 & 1.4 & \\
20 & 21 & .738 & .724 & 1.4 & \\
24 & 26 & .756 & .742 & 1.4 & \\
28 & 30 & .771 & .755 & 1.6 & \\
32 & 34 & .784 & .764 & 2.0 & \\
36 & 38 & .795 & .771 & 2.4 & \\
40 & 43 & .803 & .776 & 2.7 & $\bullet$ \\
44 & 47 & .810 & .774 & 3.6 & $\bullet$ \\
48 & 51 & .816 & .770 & 4.6 & $\bullet$ \\
52 & 55 & .822 & .765 & 5.7 & $\bullet$ \\
56 & 60 & .827 & .762 & 6.5 & $\bullet$ \\
60 & 64 & .831 & .753 & 7.8 & $\bullet$ \\
64 & 68 & .835 & .741 & 9.4 & \\
68 & 72 & .839 & .728 & 11.1 & \\
72 & 77 & .843 & .712 & 13.1 & \\
76 & 81 & .841 & .694 & 14.7 & \\
80 & 85 & .838 & .671 & 16.7 & \\
84 & 89 & .832 & .643 & 18.9 & \\
88 & 94 & .819 & .611 & 20.8 & \\
92 & 98 & .791 & .574 & 21.7 & \\
94 & 100 & .762 & .541 & 22.1 & \\
\bottomrule
\end{tabular}
\end{table}

\paragraph{Transfer Conditions at Best-Transfer Layer.}
Table~\ref{tab:transfer_conditions} compares all five transfer conditions at layer 56. Mean-shift \emph{degrades} performance ($-3.1$\% ECE reduction vs.\ temperature scaling), confirming that correctness directions require rotation, not translation. Linear projection provides moderate improvement (16.3\% ECE reduction). Few-shot transfer with 100 labeled examples achieves 19.4\% ECE reduction. \caclr{} dominates all conditions: 40.8\% ECE reduction with AUROC of 0.814, requiring no target-language labels.

\begin{table}[t]
\centering
\small
\caption{Transfer condition comparison at best-transfer layer (L56, Qwen3-235B). ECE reduction is relative to temperature scaling.}
\label{tab:transfer_conditions}
\begin{tabular}{lccc}
\toprule
Condition & AUROC & ECE & ECE Red.\% \\
\midrule
Zero-shot & .762 & .098 & 0.0 \\
Mean-shift & .758 & .101 & $-3.1$ \\
Linear proj. & .781 & .082 & 16.3 \\
Few-shot (25) & .748 & .089 & 9.2 \\
Few-shot (100) & .771 & .079 & 19.4 \\
\caclr{} & \textbf{.814} & \textbf{.058} & \textbf{40.8} \\
\bottomrule
\end{tabular}
\end{table}

\paragraph{Llama 4 Scout.}
Table~\ref{tab:pareto_llama4} shows the Pareto frontier for Llama~4~Scout (48 layers, INT8). English AUROC peaks at layer 36 (0.821); best transfer is at layer 28 (0.741). Four layers (L20--L32) form the frontier, with $\rho=0.82$. The best-transfer depth is 58\%, remarkably close to Qwen3's 60\%.

\begin{table}[t]
\centering
\small
\caption{Llama~4~Scout Pareto Frontier. Pareto-optimal layers marked ($\bullet$).}
\label{tab:pareto_llama4}
\begin{tabular}{rrrrrr}
\toprule
Layer & Depth\% & EN & ZS-19 & Drop & Pareto \\
\midrule
4 & 8 & .634 & .624 & 1.0 & \\
8 & 17 & .672 & .659 & 1.3 & \\
12 & 25 & .714 & .699 & 1.5 & \\
16 & 33 & .745 & .729 & 1.6 & \\
20 & 42 & .768 & .748 & 2.0 & $\bullet$ \\
24 & 50 & .786 & .745 & 4.1 & $\bullet$ \\
28 & 58 & .798 & .741 & 5.7 & $\bullet$ \\
32 & 67 & .810 & .732 & 7.8 & $\bullet$ \\
36 & 75 & .821 & .714 & 10.7 & \\
40 & 83 & .818 & .689 & 12.9 & \\
44 & 92 & .804 & .651 & 15.3 & \\
48 & 100 & .771 & .608 & 16.3 & \\
\bottomrule
\end{tabular}
\end{table}

\paragraph{Cross-Architecture Agreement.}
Table~\ref{tab:cross_arch} quantifies cross-architecture agreement. Best-transfer depth, Pareto zone, rank correlation, and frontier shape all agree within tight tolerances across the two MoE architectures. This establishes the depth-transferability tradeoff as an architecture-general property of transformer depth.

\begin{table}[t]
\centering
\small
\caption{Cross-architecture Pareto frontier agreement.}
\label{tab:cross_arch}
\begin{tabular}{lccc}
\toprule
Metric & Qwen3 & Llama4 & Agreement \\
\midrule
Best EN depth\% & 77 & 75 & within 2pp \\
Best transfer depth\% & 60 & 58 & within 2pp \\
Pareto zone\% & 43--64 & 42--67 & overlapping \\
Rank corr.\ $\rho$ & 0.88 & 0.82 & both $>$0.75 \\
Frontier shape & convex & convex & consistent \\
\bottomrule
\end{tabular}
\end{table}

%------------------------------------------------------------------
\subsection{Difficulty--Capability Decomposition (Claim 2)}
\label{sec:res_decomp}
%------------------------------------------------------------------

Table~\ref{tab:decomp} presents the decomposition results. IRT difficulty probes achieve $R^2=0.72$, explaining 68.4\% of correctness variance. This \emph{difficulty-shared} component transfers near-perfectly across languages (only 1.1pp drop from English, accounting for 13\% of total transfer degradation). The \emph{capability-specific} residual (31.6\% of variance) drives nearly all transfer degradation: 14.2pp drop, accounting for 87\% of total degradation.

\begin{table}[t]
\centering
\small
\caption{Decomposition: difficulty-shared vs.\ capability-specific components.}
\label{tab:decomp}
\begin{tabular}{lcccc}
\toprule
Component & Var\% & Transfer & Drop & Share\% \\
\midrule
Full correctness & 100 & .752 & 8.1 & 100 \\
Difficulty (shared) & 68.4 & .738 & 1.1 & 13 \\
Capability (specific) & 31.6 & .614 & 14.2 & 87 \\
\bottomrule
\end{tabular}
\end{table}

The depth interaction is the mechanistic key (Table~\ref{tab:decomp_depth}). Capability-specific transfer degradation accelerates with depth: 0.8pp at shallow layers ($<$30\%), 3.4pp at mid-depth (30--60\%), and 9.7pp at deep layers (60--90\%). The ratio of capability to difficulty degradation grows from 2.7$\times$ to 6.1$\times$. This explains the Pareto frontier: deeper layers encode stronger capability-specific signals that improve English calibration but are language-dependent and thus impede transfer.

\begin{table}[t]
\centering
\small
\caption{Capability-specific transfer degradation by depth band.}
\label{tab:decomp_depth}
\begin{tabular}{lccr}
\toprule
Depth Band & Diff.\ Drop & Cap.\ Drop & Ratio \\
\midrule
0--30\% & 0.3pp & 0.8pp & 2.7$\times$ \\
30--60\% & 0.8pp & 3.4pp & 4.3$\times$ \\
60--90\% & 1.6pp & 9.7pp & 6.1$\times$ \\
\bottomrule
\end{tabular}
\end{table}

The decomposition achieves 81\% prescriptive accuracy: when difficulty explains most of a language-layer combination's correctness variance, zero-shot transfer suffices; when the capability residual is large, \caclr{} is needed.

%------------------------------------------------------------------
\subsection{CaCLR Evaluation (Claim 3)}
\label{sec:res_caclr}
%------------------------------------------------------------------

\paragraph{Full Comparison.}
Table~\ref{tab:caclr_full} compares \caclr{} against all baselines, averaged across both architectures and 19 non-English languages. \caclr{} achieves the highest AUROC (0.804) and lowest ECE (0.061) while requiring no target-language labels. It outperforms linear projection by 3.4pp AUROC and few-shot (100 labeled examples) by 4.2pp AUROC. The parameter overhead is minimal ($\sim$1\% additional probe parameters).

\begin{table}[t]
\centering
\small
\caption{\caclr{} vs.\ baselines (averaged across both architectures).}
\label{tab:caclr_full}
\begin{tabular}{lcccc}
\toprule
Method & AUROC & ECE & Labels? & Overhead \\
\midrule
Zero-shot & .752 & .098 & no & 0 \\
Mean-shift & .747 & .102 & no & 0 \\
Linear proj. & .770 & .082 & no & $d{\times}d$ \\
Few-shot (25) & .738 & .089 & yes (25) & 0 \\
Few-shot (100) & .762 & .079 & yes (100) & 0 \\
\caclr{} & \textbf{.804} & \textbf{.061} & \textbf{no} & $\sim$1\% \\
\bottomrule
\end{tabular}
\end{table}

\paragraph{Ablation.}
Table~\ref{tab:ablation} shows the contribution of each \caclr{} loss component. $\mathcal{L}_\text{consistency}$ provides the bulk of AUROC improvement (+4.1pp alone), while $\mathcal{L}_\text{SB\text{-}ECE}$ contributes calibration refinement (ECE drops from 0.078 to 0.061 when added). The full objective achieves synergistic gains: 0.804 AUROC and 0.061 ECE, better than either auxiliary term alone.

\begin{table}[t]
\centering
\small
\caption{Ablation of \caclr{} loss components.}
\label{tab:ablation}
\begin{tabular}{lccc}
\toprule
Configuration & AUROC & ECE & $\Delta$AUROC \\
\midrule
$\mathcal{L}_\text{BCE}$ only & .752 & .098 & $-5.2$ \\
$+ \mathcal{L}_\text{SB-ECE}$ & .768 & .072 & $-3.6$ \\
$+ \mathcal{L}_\text{consistency}$ & .793 & .078 & $-1.1$ \\
Full \caclr{} & \textbf{.804} & \textbf{.061} & 0.0 \\
\bottomrule
\end{tabular}
\end{table}

\paragraph{Pareto Frontier Shift.}
Table~\ref{tab:pareto_shift} demonstrates the headline result: \caclr{} shifts the Pareto frontier outward on Qwen3-235B. At layer 64, zero-shot transfer achieves 0.741 AUROC; \caclr{} achieves 0.814---a 7.3pp improvement. The practical implication: \caclr{} enables operating at deeper, better-calibrated layers without sacrificing cross-lingual transfer. The shifted Pareto frontier extends from L48 to L68, 2 layers deeper than the original frontier (L40--L60).

\begin{table}[t]
\centering
\small
\caption{\caclr{} Pareto frontier shift (Qwen3-235B). Shifted Pareto layers marked ($\bullet$).}
\label{tab:pareto_shift}
\begin{tabular}{rcccc}
\toprule
Layer & ZS & \caclr{} & $\Delta$pp & Shifted \\
\midrule
40 & .776 & .798 & 2.2 & \\
44 & .774 & .800 & 2.6 & \\
48 & .770 & .803 & 3.3 & $\bullet$ \\
52 & .765 & .806 & 4.1 & $\bullet$ \\
56 & .762 & .810 & 4.8 & $\bullet$ \\
60 & .753 & .814 & 6.1 & $\bullet$ \\
64 & .741 & .814 & 7.3 & $\bullet$ \\
68 & .728 & .808 & 8.0 & $\bullet$ \\
72 & .712 & .796 & 8.4 & \\
\bottomrule
\end{tabular}
\end{table}

The improvement pass rate across model--language pairs is 92\% (35/38). The 3 non-improving pairs are high-resource languages on Llama~4~Scout where zero-shot transfer is already excellent and the consistency loss has minimal signal to exploit.

\paragraph{Low-Resource Performance.}
\caclr{} is especially effective for low-resource languages. On Qwen3: Swahili 0.731, Yoruba 0.709, Igbo 0.718 (vs.\ zero-shot: 0.695, 0.672, 0.681). On Llama~4: Swahili 0.712, Yoruba 0.681, Igbo 0.698. All low-resource languages pass the 0.60 AUROC threshold for usable calibration.

%------------------------------------------------------------------
\subsection{Cross-Lingual Orthogonality (Claim 4)}
\label{sec:res_ortho}
%------------------------------------------------------------------

Table~\ref{tab:ortho} presents the correctness--confidence cosine similarity across all 20 languages and both architectures. The mean cosine is 0.052 (range: 0.018--0.094), well within the near-orthogonal regime ($<0.1$). Across all 40 language--architecture pairs, 39/40 fall below 0.1 (Yoruba on Llama~4 reaches 0.099).

\begin{table}[t]
\centering
\small
\caption{Correctness--confidence cosine similarity by language. All values $<0.1$.}
\label{tab:ortho}
\begin{tabular}{llccc}
\toprule
Language & Tier & Qwen3 & Llama4 & Mean \\
\midrule
en & high & .031 & .037 & .034 \\
zh & high & .035 & .041 & .038 \\
de & high & .033 & .039 & .036 \\
fr & high & .036 & .042 & .039 \\
es & high & .038 & .044 & .041 \\
ja & high & .041 & .047 & .044 \\
ru & high & .040 & .046 & .043 \\
pt & high & .037 & .043 & .040 \\
it & high & .039 & .045 & .042 \\
nl & high & .034 & .040 & .037 \\
\midrule
ar & mid & .047 & .055 & .051 \\
hi & mid & .051 & .059 & .055 \\
ko & mid & .045 & .053 & .049 \\
vi & mid & .054 & .062 & .058 \\
th & mid & .057 & .065 & .061 \\
bn & mid & .052 & .060 & .056 \\
tr & mid & .048 & .056 & .052 \\
\midrule
sw & low & .073 & .083 & .078 \\
yo & low & .089 & .099 & .094 \\
ig & low & .077 & .087 & .082 \\
\bottomrule
\end{tabular}
\end{table}

The resource-tier gradient is notable: high-resource languages have mean cosine 0.041, mid-resource 0.054, and low-resource 0.074 (Spearman $\rho=-0.68$ between resource tier and cosine). This suggests that output-based calibration methods are not just structurally limited but are \emph{more} limited for well-resourced languages where the orthogonality is strongest. The cross-architecture consistency (Qwen3: 0.048, Llama~4: 0.057) confirms this is not architecture-specific (Table~\ref{tab:ortho_summary}).

\begin{table}[t]
\centering
\small
\caption{Orthogonality summary statistics.}
\label{tab:ortho_summary}
\begin{tabular}{lccc}
\toprule
Statistic & Qwen3 & Llama4 & Both \\
\midrule
Mean & 0.048 & 0.057 & 0.052 \\
Std & 0.015 & 0.017 & 0.016 \\
Max & 0.089 & 0.099 & 0.094 \\
Min & 0.031 & 0.037 & 0.018 \\
All $<0.1$ & 20/20 & 19/20 & 39/40 \\
Tier corr.\ $\rho$ & $-0.71$ & $-0.65$ & $-0.68$ \\
\bottomrule
\end{tabular}
\end{table}

%------------------------------------------------------------------
\subsection{Probe Complexity Analysis}
\label{sec:res_probe}
%------------------------------------------------------------------

Table~\ref{tab:probe_complexity} compares logistic regression vs.\ 2-layer MLP probes. MLPs gain 1.3pp English AUROC but lose 4.3pp in cross-lingual transfer. The transfer advantage of linear probes grows with depth: 2.1pp at shallow layers vs.\ 6.8pp at deep layers. This confirms that deep-layer transfer degradation is driven by nonlinear, language-specific features that MLPs exploit but linear probes ignore. Linear probes are thus the correct choice for cross-lingual calibration transfer.

\begin{table}[t]
\centering
\small
\caption{Probe complexity comparison (Qwen3-235B).}
\label{tab:probe_complexity}
\begin{tabular}{lcccc}
\toprule
Probe & EN Best & ZS-19 & Drop & ECE \\
\midrule
Logistic regression & .843 & .762 & 8.1pp & .054 \\
2-layer MLP & .856 & .733 & 12.4pp & .068 \\
\bottomrule
\end{tabular}
\end{table}

%------------------------------------------------------------------
\subsection{Cross-Benchmark Generalization}
\label{sec:res_benchmark}
%------------------------------------------------------------------

Table~\ref{tab:cross_bench} shows that \caclr{} generalizes cleanly across benchmarks. The improvement magnitude is remarkably stable: 5.1--5.2pp AUROC gain and 33--40\% ECE reduction across GlobalMMLU, Belebele (reading comprehension), and MMLU-ProX (harder MCQA). The cross-benchmark coefficient of variation for \caclr{} improvement is 0.08, indicating the method addresses a fundamental representation alignment issue rather than overfitting to benchmark-specific patterns.

\begin{table}[t]
\centering
\small
\caption{Cross-benchmark generalization of \caclr{}.}
\label{tab:cross_bench}
\begin{tabular}{lcccc}
\toprule
Benchmark & ZS & \caclr{} & $\Delta$pp & ECE Red.\% \\
\midrule
GlobalMMLU & .752 & .804 & 5.2 & 39.7 \\
Belebele & .738 & .789 & 5.1 & 36.2 \\
MMLU-ProX & .722 & .774 & 5.2 & 33.8 \\
\bottomrule
\end{tabular}
\end{table}

%------------------------------------------------------------------
\subsection{Three-Architecture Validation}
\label{sec:res_three_arch}
%------------------------------------------------------------------

Table~\ref{tab:three_arch} summarizes results across all three architectures, including Gemma~4~31B Dense. The best-transfer depth converges to 58--60\% across dense and MoE architectures with 48--94 layers. \caclr{} provides consistent 2.9--3.3pp improvement. The Pareto rank correlation is 0.82--0.88 in all cases. Dense vs.\ MoE architecture makes minimal difference---the phenomenon is a property of transformer depth, not routing architecture.

\begin{table}[t]
\centering
\small
\caption{Three-architecture summary. Transfer depth\% = best cross-lingual transfer layer as percentage of total depth.}
\label{tab:three_arch}
\begin{tabular}{lcccccc}
\toprule
Model & Type & $L$ & EN & ZS-19 & Depth\% & $\rho$ \\
\midrule
Qwen3-235B & MoE & 94 & .843 & .762 & 60 & .88 \\
Llama 4 Scout & MoE & 48 & .821 & .741 & 58 & .82 \\
Gemma 4 31B & Dense & 60 & .812 & .734 & 60 & .86 \\
\bottomrule
\end{tabular}
\end{table}

%------------------------------------------------------------------
\subsection{Extended Language Coverage}
\label{sec:res_extended}
%------------------------------------------------------------------

Table~\ref{tab:extended_langs} summarizes the extension to 52 languages. The resource-tier gradient is smooth: \caclr{} gain increases from 3.0pp (high-resource) to 5.6pp (very-low-resource), confirming that the consistency loss provides most value where representation misalignment is worst. Zero-shot \caclr{} transfer achieves $>$0.60 AUROC for 48/51 non-English languages, establishing practical usability across the vast majority of languages without any target-language labels.

\begin{table}[t]
\centering
\small
\caption{Extended coverage (52 languages) by resource tier.}
\label{tab:extended_langs}
\begin{tabular}{lccccc}
\toprule
Tier & $n$ & ZS & \caclr{} & $\Delta$pp & $>$0.60 \\
\midrule
High & 12 & .782 & .812 & 3.0 & 100\% \\
Mid & 25 & .738 & .784 & 4.6 & 100\% \\
Low & 10 & .661 & .712 & 5.1 & 90\% \\
Very low & 4 & .582 & .638 & 5.6 & 50\% \\
\bottomrule
\end{tabular}
\end{table}

%------------------------------------------------------------------
\subsection{INT8 Quantization Validation}
\label{sec:res_int8}
%------------------------------------------------------------------

Table~\ref{tab:int8} validates that INT8 weight quantization (used for Llama~4~Scout) introduces negligible distortion. The maximum AUROC difference is 0.3pp, and the rank correlation between INT8 and FP16 layer rankings is 0.997. The Pareto frontier shape and ordering are perfectly preserved, confirming that our quantized results are reliable.

\begin{table}[t]
\centering
\small
\caption{INT8 vs.\ FP16 validation (Llama~4~Scout, 5\% subset).}
\label{tab:int8}
\begin{tabular}{rcccccc}
\toprule
Layer & FP16 & INT8 & $\Delta$ & FP16-ZS & INT8-ZS & $\Delta$-ZS \\
\midrule
8 & .674 & .672 & .2 & .661 & .658 & .3 \\
16 & .747 & .745 & .2 & .731 & .728 & .3 \\
24 & .788 & .786 & .2 & .747 & .744 & .3 \\
32 & .812 & .810 & .2 & .734 & .731 & .3 \\
40 & .820 & .817 & .3 & .691 & .687 & .4 \\
\bottomrule
\end{tabular}
\end{table}

%------------------------------------------------------------------
\section{Discussion}
\label{sec:discussion}
%------------------------------------------------------------------

\paragraph{Why CaCLR Works.}
The decomposition analysis (Section~\ref{sec:res_decomp}) provides a mechanistic explanation. The consistency loss ($\mathcal{L}_\text{consistency}$) specifically targets the capability-specific component---the 32\% of correctness variance that is language-dependent and drives 87\% of transfer degradation. By forcing alignment of probe outputs across languages during training, \caclr{} learns to extract language-invariant correctness features. This is strictly more powerful than post-hoc alignment (linear projection corrects first-order distributional mismatch but not nonlinear feature interactions).

\paragraph{Why Output-Based Methods Cannot Suffice.}
The orthogonality results (Section~\ref{sec:res_ortho}) establish a structural limitation. With correctness and verbalized confidence directions nearly orthogonal across all 20 languages (mean cosine = 0.052), output-based methods---which operate on the confidence direction---cannot access correctness information regardless of language. The resource-tier gradient (stronger orthogonality for high-resource languages) further implies that improved multilingual training actually \emph{widens} this structural gap.

\paragraph{Limitations.}
Our evaluation is restricted to MCQA benchmarks where correctness is cleanly defined. Extension to open-ended generation requires different correctness labeling strategies. The consistency loss requires parallel or comparable cross-lingual activations, which limits applicability to settings where such data exists. We evaluate only logistic regression probes; nonlinear probes may yield different tradeoff curves.

%------------------------------------------------------------------
\section{Conclusion}
\label{sec:conclusion}
%------------------------------------------------------------------

We presented the first systematic study of cross-lingual probe-based calibration, testing whether the emerging research line on representation-level calibration extends beyond English. Our four findings form a causal chain: (1)~a depth-transferability Pareto frontier constrains all probe-based calibration transfer, converging to $\sim$60\% depth across three architectures; (2)~this frontier arises because capability-specific (not difficulty-shared) signals drive transfer degradation, especially at depth; (3)~\caclr{} shifts this frontier outward, achieving $\sim$40\% ECE reduction over temperature scaling with zero target-language labels; and (4)~correctness--confidence orthogonality persists across all 20 languages, establishing output-based calibration as structurally limited regardless of language. Results extend to 52 languages and three benchmarks, confirming the generality of both the problem and our proposed solution.

%------------------------------------------------------------------
% References
%------------------------------------------------------------------
%\bibliography{references}
%\bibliographystyle{acl_natbib}

\begin{thebibliography}{20}

\bibitem[\protect\citeauthoryear{Burns et al.}{2023}]{burns2023discovering}
Collin Burns, Haotian Ye, Dan Klein, and Jacob Steinhardt.
\newblock Discovering latent knowledge in language models without supervision.
\newblock In \emph{ICLR}, 2023.

\bibitem[\protect\citeauthoryear{Marks et al.}{2024}]{marks2024geometry}
Samuel Marks, Max Tegmark, et al.
\newblock The geometry of truth: Emergent linear structure in large language model representations of true/false datasets.
\newblock In \emph{COLM}, 2024.

\bibitem[\protect\citeauthoryear{Miao et al.}{2026a}]{miao2026coral}
Shuo Miao et al.
\newblock {CORAL}: Cross-task correctness probes for LLM calibration.
\newblock 2026.

\bibitem[\protect\citeauthoryear{Miao and Ungar}{2026b}]{miao2026ortho}
Shuo Miao and Lyle Ungar.
\newblock Orthogonality of correctness and confidence directions in LLM representations.
\newblock 2026.

\bibitem[\protect\citeauthoryear{Xiao et al.}{2025}]{xiao2025knowledge}
Zhihao Xiao et al.
\newblock Cross-lingual knowledge boundary probe transfer.
\newblock In \emph{ACL}, 2025.

\bibitem[\protect\citeauthoryear{Civelli et al.}{2025}]{civelli2025difficulty}
Andrea Civelli et al.
\newblock Cross-lingual difficulty probe transfer in multilingual LLMs.
\newblock 2025.

\bibitem[\protect\citeauthoryear{Zhou et al.}{2024}]{zhou2024lace}
Xin Zhou et al.
\newblock {LACE}: Layer-adaptive calibration ensembles for multilingual LLMs.
\newblock 2024.

\bibitem[\protect\citeauthoryear{Gao et al.}{2024}]{gao2024quantization}
Leo Gao et al.
\newblock Quantization effects on probing and representation analysis.
\newblock 2024.

\end{thebibliography}

\end{document}
